\chapter{Допуск геометрически растущего общего носителя}
\label{ch:v10-quantum-rg-common-carrier-admission}

\section{Переход от статического масштаба к истории геометрии}

Финальный запрет Тома IX показал, что абсолютный масштаб нельзя получить
переименованием радиуса, температуры, частоты часов или масштаба отсечения.
Настоящий гейт проверяет другую возможность: фундаментальным объектом является
не одно число $E_*$, а растущая геометрическая история, а локальный
спектральный атлас задаёт устойчивый безразмерный рисунок на этой истории.

Пусть $N(\tau)>0$ --- число элементарных пространственных ячеек на срезе с
безразмерным порядковым параметром $\tau$. Определим относительный масштаб
пространства
\begin{equation}
 a(\tau)=\left(\frac{N(\tau)}{N_0}\right)^{1/3}.
 \label{eq:v10-cell-scale-factor}
\end{equation}
Параметр $\tau$ пока не объявляется физическим временем. Он нумерует историю
роста и сохраняет запрет Тома VIII на автоматическое происхождение секунды.

\section{Общий носитель}

Минимальная архитектура объединяет четыре слоя:
\begin{equation}
 \mathcal H_X=
 \ell^2(\mathbb N_0)
 \widehat\otimes \mathcal K_{43}
 \widehat\otimes \mathcal H_{\rm KMS}^{(6)}
 \widehat\otimes \mathcal F_{\rm loc}.
 \label{eq:v10-common-growth-carrier}
\end{equation}
Первый слой хранит порядок рождения ячеек, второй наследует автономный
конвейер Тома VIII, третий несёт концевое KMS-завершение Тома IX, четвёртый
несёт локальные квантовые поля и безразмерный оператор Дирака
$\widehat D(g_1,\ldots,g_r)$.

Матрица инцидентности трёх основных типов данных выбрана как
\begin{equation}
 C_X=
 \begin{pmatrix}
 1&1&0\\
 0&1&1\\
 1&0&1
 \end{pmatrix},
 \qquad \rank C_X=3,
 \qquad \det C_X=2.
 \label{eq:v10-carrier-incidence}
\end{equation}
Поэтому история, ячейка и локальный полевой сектор не сводятся друг к другу
простым переименованием.

\section{Геометрический ход и космологическая кривизна}

Старая вакуумная формула проекта допускает новое точное разложение. В
планковских единицах введём условную амплитуду роста
\begin{equation}
 h_{\rm vac}=\frac{e^{-S_{\rm vac}}}{\sqrt{8\pi}}.
 \label{eq:v10-vacuum-growth-amplitude}
\end{equation}
Тогда история постоянного относительного роста имеет вид
\begin{equation}
 N(\tau)=N_0e^{3h_{\rm vac}\tau},
 \qquad
 a(\tau)=e^{h_{\rm vac}\tau},
 \label{eq:v10-cell-growth-history}
\end{equation}
и точно выполняются
\begin{equation}
 \frac{1}{a}\frac{da}{d\tau}
 =\frac{1}{3N}\frac{dN}{d\tau}
 =h_{\rm vac}.
 \label{eq:v10-growth-rate-identity}
\end{equation}

Определённая через квадрат темпа кривизна равна
\begin{equation}
 \Lambda_{\rm growth}:=3h_{\rm vac}^2
 =\frac{3}{8\pi}e^{-2S_{\rm vac}}.
 \label{eq:v10-growth-cosmological-identity}
\end{equation}
Правая часть совпадает с ранней формулой космологической постоянной проекта.
Теперь она имеет строгую алгебраическую интерпретацию как квадрат условной
амплитуды геометрического роста. Сам закон
$h_{\rm vac}=e^{-S_{\rm vac}}/\sqrt{8\pi}$ пока не выведен из меры переходов.

\section{Геометрическая координата ренормгруппового хода}

Естественная безразмерная координата масштаба определяется самой историей:
\begin{equation}
 \zeta:=\log a=\frac13\log\frac{N}{N_0}.
 \label{eq:v10-geometric-rg-coordinate}
\end{equation}
Квантовый ход локальных связей может быть записан без внешней размерной
шкалы как
\begin{equation}
 \beta_i(g):=\frac{dg_i}{d\zeta}.
 \label{eq:v10-geometric-beta-definition}
\end{equation}
Это определение создаёт общий носитель для геометрического роста и
ренормгрупповой динамики, но не утверждает ненулевость $\beta_i$.

Локальный спектральный атлас факторизуется:
\begin{equation}
 D(\zeta)=E_{\rm cell}(\zeta)\widehat D(g(\zeta)),
 \qquad
 m_i(\zeta)=E_{\rm cell}(\zeta)\mu_i(g(\zeta)).
 \label{eq:v10-local-atlas-factorization}
\end{equation}
При общем множителе отношения $m_i/m_j=\mu_i/\mu_j$ не зависят от
$E_{\rm cell}$. Поэтому геометрический атлас может быть строгим рисунком даже
до происхождения абсолютного масштаба.

\section{Две неэквивалентные ветви}

В ветви добавления ячеек их локальный объём и локальный энергетический
масштаб постоянны:
\begin{equation}
 V(\tau)=v_*N(\tau),
 \qquad v_*=\mathrm{const},
 \qquad E_{\rm cell}=E_0.
 \label{eq:v10-fixed-cell-growth-branch}
\end{equation}
Расширение означает рост числа локальных ячеек, а не растяжение каждой из них.

В ветви совместного растяжения
\begin{equation}
 E_{\rm cell}(\tau)=\frac{E_0}{a(\tau)},
 \qquad
 a(\tau)E_{\rm cell}(\tau)=E_0.
 \label{eq:v10-codilation-branch}
\end{equation}
Все локальные линейки и массы изменяются общим множителем. Без отдельного
безразмерного отношения эта ветвь является конформным переописанием и не
снимает масштабный запрет.

\begin{theorem}[Допуск растущего носителя]
Носитель \eqref{eq:v10-common-growth-carrier} совместно размещает порядковую
историю, автономную поставку ячеек, концевой KMS-сектор и локальный
квантово-полевой атлас. Формулы
\eqref{eq:v10-cell-scale-factor}--\eqref{eq:v10-growth-cosmological-identity}
точно согласуют число ячеек, масштабный фактор, относительный темп роста и
раннюю космологическую формулу. Все семь архитектурных условий выполнены.
\end{theorem}

\begin{nogocriterion}[Граница первого гейта]
Допуск носителя не выводит амплитуду рождения $h_{\rm vac}$, элементарный
объём $v_*$, физическую длительность одного шага и локальный масштаб
$E_{\rm cell}$. Нельзя объявлять формулу
$\Lambda_{\rm growth}=3h_{\rm vac}^2$ физическим выводом космологической
постоянной до построения нормированной меры переходов $N\to N+1$.
\end{nogocriterion}

\begin{protocol}[Статус первого гейта Тома X]\begin{enumerate}
\item общий геометрически растущий носитель: \textbf{допущен};
\item архитектурный счёт: \textbf{$7/7$};
\item физическое происхождение закона роста, такта и локального масштаба:
\textbf{$0/3$};
\item статический масштаб Тома IX: \textbf{не подставляется};
\item следующий гейт: \textbf{происхождение ненулевого квантового хода и
следовой аномалии относительно координаты $\zeta$}.
\end{enumerate}\end{protocol}

Машинный сертификат сохранён в
\path{s2t_v10_quantum_rg_common_carrier_admission_gate_results.json}.

\begin{thebibliography}{9}
\bibitem{v10-rideout-sorkin}
D.~Rideout, R.~Sorkin,
\emph{A Classical Sequential Growth Dynamics for Causal Sets},
Physical Review D 61 (2000), 024002.

\bibitem{v10-everpresent-lambda}
M.~Ahmed, S.~Dodelson, P.~Greene, R.~Sorkin,
\emph{Everpresent $\Lambda$},
Physical Review D 69 (2004), 103523.
\end{thebibliography}